% Transcription — fonds Alexandre Grothendieck, shelfmark 111, batch 1
% (pages 1-16, the whole folder). Established from the facsimile
% archives/batches/111/batch-01.pdf (a mirror of
% https://grothendieck.umontpellier.fr/111.pdf), per the skill
% transcribe-grothendieck. The batch file carries no cover sheet: PDF page and
% archivists' page coincide. All sixteen pages are transcribed; none is
% skipped.
%
% What the folder is. A photocopy of a sixteen-page typescript, numbered 1 to
% 16 by hand at the head of each page, typed on a machine without Greek
% letters or arrows (some inked in afterwards), with many typewriter x-outs and
% typed interlinear insertions, and then gone over by hand (the hand's marks
% are reproduced with the typing; see « Relation to folder 162-5 »). The typed text is headed « Tapis de Quillen. » on
% page 1, so that part of the inventory's title is on the page; the bracketed
% « [notes du 10/09/1968] » is the archivists' transcription of a handwritten
% « notes du 10.9.68 », underlined, at the top right of page 1. The date is
% thus the hand's, not the archivists' inference, and it is given where it
% stands; \dating{} carries the inventory's « 1968 » verbatim.
%
% Content. Notes on what Quillen told or showed the writer, in three typed
% parts. I (pp. 1-5): relations between categories and semi-simplicial sets —
% the nerve S : (Cat) -> (Ssimpl), local systems on SC as functors inverting
% arrows, their cohomology as derived limits or as cohomology of the topos
% C^, the functor T(X) = Delta/X and Quillen's theorem that S and T induce
% quasi-inverse equivalences of the homotopy categories, the canonical
% homotopism between C and C°, the sheaf-theoretic interpretation of the
% cohomology of a semi-simplicial set with local coefficients; then a
% generalisation proposed by the writer to morphisms of « localement
% homotopiquement triviaux » topoi, with the equivalence of a) homotopism,
% b) equivalence D^b_lc(X') -> D^b_lc(X), a') equivalence on abelian local
% systems plus isomorphism on H^i, and a sketch of proof. II (pp. 5-10):
% bicategories as category objects in (Cat) and bi-semi-simplicial sets,
% 2-categories, (Cat)-groups and Dold-Kan (« le tapis de Dold »): abelian
% (Cat)-groups = chain complexes of length 1 (« ma construction chérie de la
% catégorie de Picard »), abelian (2-Cat)-groups = complexes of length 2, the
% expectation that abelian groups in (n-Cat) correspond to complexes of length
% n; Gr- and Grab-n-categories as truncations of H-spaces; universal
% Gr-n-categories of a monoidal category and higher K_n(C). III (pp. 11-16):
% « point de vue motivique » in cobordism — Quillen's universal pair
% (F., F^•) on differentiable manifolds under homotopy, satisfying base change
% for transversal squares, realised by the category B with Hom_B(X, Y) =
% B(X x Y) (unoriented bordism), a proof of the universal property; the
% structure B.(X) = H^•(X) ⊗ Lambda over Lambda = B.(pt); the Karoubian
% envelope, inverting Lambda(-1), modules of finite presentation over Lambda;
% the analogy with motives for rational equivalence; a second universal
% category built on S(X) ⊂ H^•(X), the classes of submanifolds, via Thom's
% theorem in terms of Steenrod operations; the affine group scheme G over F_2
% dual to the Steenrod algebra as a « groupe de Galois motivique », strongly
% non-reductive; Quillen's generalisation to oriented, non-proper settings.
%
% Attribution — decided not to assert. The typed text is in the first person
% and says « ma lettre à Verdier » (p. 8), « "ma" construction chérie de la
% "catégorie de Picard" » (p. 7), « mon interprétation chérie » (p. 9); it
% reports what « Quillen prouve », and twice « Demander à Quillen ». Nothing
% on the pages signs it, and the edition does not name its author. The hand
% that numbered the pages, dated page 1 and annotated is likewise not
% identified here; it writes in the manner of the other hand-annotated
% typescripts of the fonds, but that is not evidence the file relies on.
%
% Layers, and how they are kept apart. The typed text is set as running text,
% typed underlinings as \emph{}; typed interlinear insertions are typed text
% and are set in place without apparatus; typewriter x-outs are not
% reproduced (they are the typing), except where an x-ed passage carries a
% thought, which is then given in a \note{}. Everything handwritten is
% apparatus: handwritten insertions into the typed line are \add{} — following
% folders 55 and 103, and with the consequence that this file uses \add{} for
% nothing else — typed words struck by hand are \struck{}, and handwritten
% notes that stand apart from the line (margins, heads and feet of pages,
% interlinear sentences) are \marginal{}. Here \marginal{} marks the
% handwritten layer and does not attribute it. Where the hand only corrects
% notation (superscripts turned into subscripts on page 5, tildes and
% asterisks inked in), the corrected form is set and a note says so once.
%
% Proofs. The typescript is notes, not a treatise; its proof sketches are a
% few lines each (p. 4, the universal property on pp. 12-13) and the hand
% works inside them (the labels of the square on p. 13), so they are given
% whole rather than summarised.
%
% The typescript's slips are left as typed and not flagged one by one:
% « fancteurs », « dérigés »,
% « coéfficiens », « boefficients », « catétorie », « homomorhisme »,
% « équivelente », « strcture », « thorème », « corespond », « ppouve »,
% « prédédente », « poihts », « néeessairement », « redactor demerdetur »
% (sic), « un triple » for a quintuple, and « C », « C' » on page 4 where the
% sentence is about topoi X, X'.
%
% Relation to folder 162-5. Folder 162-5 holds two more witnesses of the same
% typed master (same line breaks, same x-outs): copy B (162-5 pp. 32-40, 41,
% 43-48; typescript page N = 162-5 p. 31+N for N <= 9, p. 41 for N = 10,
% p. 32+N for N = 11-16), the typed top copy annotated in blue ink; and copy A
% (162-5 pp. 2-17, typescript page N = p. N+1), a black photocopy of B taken
% after the ink, carrying further notes of its own. The handwritten layer of
% 111 is B's ink, photocopied, plus the notes only A shares: « (compactes ?) »
% p. 11, « NB H*(X) (= H*(X, F_2)) » p. 13, « i.e. l'enveloppe "Karoubienne" »
% p. 13, « = karoubienne + additive » p. 14, the Λ filled into the blank of
% p. 14 (« libres de type fini sur Λ »), « cf exposé Karoubi à Bourbaki »
% p. 16. On the images these A-only notes have, on 111 and on A, the same
% strokes in the same places (same breaks, same slant, same wavy underline
% under « quasi-abélienne »): they are one set of marks reproduced twice,
% not written again on 111, and on 111 nothing in their tone separates them
% from the photocopied ink (the « au crayon » of the first pass is not
% supported; they are grey like the rest). 111 is nonetheless not simply a
% photocopy of A as A was scanned: on several leaves (pp. 7, 13) 111 keeps
% text at the right edge that A's leaf has lost (« déter », « prouve »,
% « ,et »). What the images support: 111 and A reproduce one annotated state
% of B plus A's notes; whether 111 was taken from A before A's edge was lost
% or from a common intermediate, they do not say. No mark was found on 111
% that is on neither A nor B.
%
% Revision. Read against copy B (the clearest witness of the ink), with 111's
% own page for every change: a reading is corrected only where 111's image
% bears it out; where the photocopy loses what B shows, \uncertain/\ill stay
% and a \note says what copy B reads. Corrected: « semi-simpliciaux »,
% « semi-simplicial » (typed overstrikes, p. 1); « des H^i( du topos »
% (p. 1); « coéfficients ensemblistes sur C (déf = » (p. 1); « i.e.
% canoniquement isomorphe » (p. 2); « covariant », doubtful (p. 3); « X et
% X' » as ink over struck « X (resp. X') », « soit … trivial », ink « ≥ 1 »,
% « de X' » struck (p. 3); the interlinear of p. 7 (« Quand on simplifie mod
% "Grab"-équivalence … ») and the margin of p. 8 (« non, on a un foncteur pl.
% fidèle (?) mais pas ess. surjectif »); « objets connexes de la catégorie
% homotopique ponctuée » (p. 8); the struck half of p. 10's paragraph; p. 12
% « F• : C → B » (bare C); p. 13 α_i without asterisk, the label Γ = (p,
% id_X), the generators x_i, « (il est douteux qu'il soit » as the hand's
% parenthesis around « compatible avec multiplication » (not struck), a)
% struck, « [Quillen ignore si » over struck « Le » before b); p. 14 « de
% tous », « à beaucoup », « convient », « lim »; p. 15 the struck word before
% H^1(S^1); p. 16 the slanted marginal (the same as B's « oui, car son anneau
% affine … stables par Δ »), « Quillen » with its t struck by hand, f^•.
%
% Doubtful readings, not settled: the hand's product sign against the typed
% coproduct on page 3; « covariant » on page 3 (con + x-outs + variant); the
% word after « bien par » in the foot note of page 3; the left-margin block of
% page 7, cancelled; words of the interlinear of page 7 (« simplifie »,
% « Grab », « seuls invariants », H_0, H_1), which the photocopy pales; the
% struck second half of the handwritten paragraph of page 10; the generators'
% condition on page 13 (« entier ≥ 1 », « ≠ de la forme 2^ν − 1 »), pale and
% clipped; the end of the slanted marginal of page 16; « convient » on page
% 14 (an overstrike).
%
% Pass: Opus 5 (claude-opus-5), 2026-09-14 — first pass, unchecked against the pages by a human.
% Revised: Opus 5.5 (claude-opus-5-5), 2026-09-22 — read against copy B, folder 162-5 pp. 32–48.
\documentclass[11pt,a4paper]{article}
% Le préambule commun à toutes les transcriptions du fonds Grothendieck.
%
% Il tient en une idée : **l'incertitude se compose, elle ne se résout pas.**
% Une transcription qui lisse un mot douteux a détruit la seule chose pour
% laquelle on la faisait — un lecteur doit pouvoir savoir, à chaque endroit,
% ce qui était sur le feuillet et ce qui a été ajouté. Les six macros qui
% suivent sont donc l'appareil critique, et non de la décoration.
%
% Les mêmes commandes sont reconnues par `scripts/render.mjs`, qui produit la
% vue de lecture du site. Une macro ajoutée ici doit l'être aussi là-bas,
% faute de quoi le rendu échouera bruyamment — ce qui est voulu : un rendu
% silencieusement faux est pire qu'un rendu absent.

\ProvidesPackage{grothendieck}[2026/08/08 Transcriptions du fonds Grothendieck]

\usepackage{amsmath,amssymb,amsthm}
\usepackage{mathrsfs}
\usepackage{stmaryrd}
% Les diagrammes commutatifs sont partout dans ce fonds : c'est la langue
% même de la géométrie algébrique de Grothendieck, et une transcription qui
% les rendrait en prose ne transcrirait rien.
\usepackage{tikz-cd}
% Français par défaut — c'est la langue des feuillets — mais l'anglais est
% chargé aussi : la traduction et le résumé sont des documents anglais, et la
% typographie française (espaces avant les deux-points) y serait une faute.
% Ces documents basculent par \selectlanguage{english} en tête de corps.
\usepackage[english,french]{babel}
\usepackage{fontspec}
\usepackage{geometry}
\geometry{a4paper,margin=25mm,marginparwidth=28mm}
\usepackage{marginnote}
\usepackage{xcolor}
% ulem, not soul. `soul`'s \st and \ul are built on a character-by-character
% scan that XeLaTeX's font machinery breaks: the document fails with an
% undefined control sequence rather than degrading. ulem does the same two jobs
% and is Unicode-engine safe. `normalem` stops it hijacking \emph, which the
% transcriptions use for Grothendieck's own emphasis.
\usepackage[normalem]{ulem}
\usepackage{hyperref}
\hypersetup{colorlinks=true,linkcolor=black,urlcolor=black,pdfborder={0 0 0}}

% --- Métadonnées du lot -----------------------------------------------------
% Reprises telles quelles par le script de rendu, qui les affiche en tête de
% la vue de lecture. Elles fixent ce que le fichier prétend couvrir : sans
% elles, une transcription flotte sans point d'attache dans un fonds de seize
% mille feuillets.

\newcommand{\folder}[1]{\gdef\grothFolder{#1}}
\newcommand{\batch}[1]{\gdef\grothBatch{#1}}
\newcommand{\leaves}[2]{\gdef\grothFirst{#1}\gdef\grothLast{#2}}
\newcommand{\foldertitle}[1]{\gdef\grothFoldertitle{#1}}
\gdef\grothFolder{?}\gdef\grothBatch{?}\gdef\grothFirst{?}\gdef\grothLast{?}\gdef\grothFoldertitle{}

% --- Le résumé --------------------------------------------------------------
%
% Il ouvre la lecture modernisée, sous le titre « Résumé », et il est composé
% en petit. Ce n'est pas un ornement : c'est le seul passage du document qui
% ne soit pas des mathématiques, et un lecteur doit pouvoir le reconnaître
% avant de l'avoir lu — puis le sauter s'il connaît déjà le sujet, ou s'y
% arrêter s'il ne le connaît pas. Le filet qui le suit marque la reprise du
% corps.

\newenvironment{resume}
  {\small\setlength{\parskip}{0.45em}}
  {\par\medskip\hrule\medskip}

% --- Mots-clés ---------------------------------------------------------------
%
% En anglais, en clôture du résumé de la lecture modernisée : le vocabulaire
% moderne sous lequel chercher ce que les feuillets construisent. Le manifeste
% du site (`npm run manifest`) les extrait pour étiqueter la cote entière —
% cette ligne est la source unique des tags, il n'y a pas de fichier à côté.
\newcommand{\keywords}[1]{\par\smallskip\noindent
  {\footnotesize\textsc{Keywords} --- \textit{#1}}\par}

% --- L'appareil critique ----------------------------------------------------

% Le numéro de feuillet, en marge. C'est l'ancre qui permet de régler un
% désaccord sur une formule en regardant *un* feuillet plutôt qu'une chemise
% de six cents. Le site s'en sert aussi pour tourner les pages du fac-similé
% quand on fait défiler la transcription.
\newcommand{\leaf}[1]{%
  \marginnote{\footnotesize\color{gray!70}\textsf{#1}}%
  \label{leaf:#1}\ignorespaces}

% Illisible : on ne devine pas. Un mot inventé qui se lit comme les autres est
% le pire résultat possible.
\newcommand{\ill}{\textcolor{red!60!black}{[\,\ldots\,]}}

% Lecture douteuse : le mot est donné, et signalé comme incertain.
\newcommand{\uncertain}[1]{\uline{#1}}

% Ajout de l'éditeur — une lettre manquante, un mot sous-entendu.
\newcommand{\add}[1]{\textcolor{blue!50!black}{[#1]}}

% Biffé par Grothendieck lui-même. Ce qu'il a rayé fait partie du document :
% une rature dit souvent où la pensée a changé de direction.
\newcommand{\struck}[1]{\sout{#1}}

% Note du transcripteur, sur l'état matériel du feuillet.
\newcommand{\note}[1]{\par\smallskip{\small\color{gray!60!black}\textit{[#1]}}\par\smallskip}

% Note marginale de Grothendieck, distincte du corps du texte.
\newcommand{\marginal}[1]{\par\smallskip\noindent\hspace{1em}%
  {\small\color{gray!60!black}#1}\par\smallskip}

% --- Titre ------------------------------------------------------------------

\AtBeginDocument{%
  \begin{center}
    {\large\bfseries Fonds Alexandre Grothendieck --- cote n\textsuperscript{o}~\grothFolder}\\[2pt]
    {\itshape\grothFoldertitle}\\[4pt]
    {\small feuillets \grothFirst--\grothLast\ (lot \grothBatch)}\\[2pt]
    {\footnotesize\color{gray!60!black}Transcription établie d'après le fac-similé numérisé
     par l'Université de Montpellier.\\
     Les lectures incertaines sont soulignées, les illisibles notées [\,\ldots\,],
     les ajouts entre crochets.}
  \end{center}
  \vspace{1em}}

\endinput


\folder{111}
\batch{1}
\pages{1}{16}
\dating{1968}
\watermark{Édition de démonstration}
\foldertitle{Tapis de Quillen [notes du 10/09/1968] : copies de tapuscrit annoté (1968)}

\begin{document}

\page{1}

\marginal{notes du \emph{10.9.68}}\note{à la main, en haut à droite de la
page, la date soulignée. Le dossier est la copie d'un tapuscrit de seize
pages, numérotées 1 à 16 à la main ; le texte tapé est donné en clair, les
interventions manuscrites en appareil : les insertions manuscrites dans la ligne sont données entre crochets, les mots tapés biffés à la main sont barrés, les notes manuscrites détachées du texte sont données à part, sans que cette présentation attribue la main. Les biffures à la machine ne sont pas reproduites}

\section*{Tapis de Quillen.}

\subsection*{I \quad Relations entre catégories et ensembles semi-simpliciaux.}

A toute catégorie $C$, on associe un ensemble semi-simplicial $S(C)$,
trouvant ainsi un foncteur pleinement fidèle
\[
S : (\mathrm{Cat}) \longrightarrow (\mathrm{Ssimpl}) .
\]
Les systêmes locaux d'ensemble sur $SC$ correspondent aux fancteurs sur $C$
qui transforment toute flêche en isomorphisme (i.e.\ qui se factorisent par le
groupoïde associé à $C$). Les $H^{i}$ sur $SC$ d'un tel systême local ($H^{0}$
pour ens, $H^{1}$ pour groupes, $H^{i}$ quelconques pour groupes abéliens)
s'interprètent en termes des foncteurs $\varprojlim^{(i)}$ dérigés de
$\varprojlim$, ou si on préfère, des $H^{i}$( du \emph{topos} $\widehat{C}$. On voit
ainsi à quelle condition un foncteur $C \to C'$ induit un homotopisme $SC \to
SC'$ : en vertu du critère cohomologique de Artin-Mazur, il f et s que pour
tout systême de coéfficients $F'$ sur $C'$, l'homomorphisme naturel
$\varprojlim_{C'}^{(i)} F' \to \varprojlim_{C}^{(i)} F$ soit un isomorphisme
(pour les $i$ pour lesquels cela a un sens).

A $C$ on peut assocér le topos $\widetilde{C}$, qui varie de façon
\emph{covariante} avec $C$. (NB le foncteur $C \mapsto \widetilde{C}$ n'a plus
rien de pleinement fidèle, semble-til ??)\note{les tildes sont tracés à la
main sur les $C$ de ce paragraphe et du suivant ; le paragraphe précédent
écrit $\widehat{C}$, d'un chapeau pâle, et ouvre devant « du \emph{topos} » une parenthèse qui n'est pas refermée}

Les systêmes de coéfficients ensemblistes sur $C$ (déf = les foncteurs
$C^{\circ} \to \mathrm{Ens}$ transformant isomorphismes en isomorphismes)
correspondent aux faisceaux localement constants i.e.\ les objets localement
constants de $\widetilde{C}$, définis intrinsèquement en termes de
$\widetilde{C}$. Ainsi, le fait pour un foncteur $F : C \to C'$ d'induire un
homotopisme $S(C) \to S(C')$ ne dépend que du morphisme de topos
$\widetilde{F} : \widetilde{C} \to \widetilde{C}'$ induit, et signifie que
pour tout faisceau localement constant $F'$ sur $C'$ i.e.\ sur
$\widetilde{C}'$, les applications induites
\[
H^{i}(\widetilde{C}', F') \longrightarrow H^{i}(\widetilde{C},
\widetilde{F}^{*}(F'))
\]
sont des isomorphismes (pour les $i$ pour lesquels cela a un sens).

\page{2}

On a aussi un foncteur évident
\[
T : (\mathrm{Ssimpl}) \longrightarrow (\mathrm{Cat}) ,
\]
en associant à tout ensemble semi-simplicial $X$ la catégorie $T(X) =
\Delta/X$ des simplexes sur $X$, dont l'ensemble des objets est la réunion
disjointe des $X_{n}$ ... (c'est une catégorie fibrée sur la catégorie
$\Delta$ des simplexes types, à fibres les catégories discrètes définies par
les $X_{n}$). Ceci posé, Quillen prouve que pour tout $X$, $ST(X)$ est
isomorphe canoniquement à $X$ dans la catégorie homotopique construite avec
$(\mathrm{Ssimpl})$, et que pour toute $C$, la catégorie $TS(C)$ est
canoniquement « homotopiquement équivalente à $C$ » i.e.\ canoniquement isomorphe à $C$ dans
la catégorie quotient de $(\mathrm{Cat})$ obtenue en inversant les foncteurs
qui sont des homotopismes. Ces isomorphismes sont fonctoriels en $X$. Il en
résulte formellement qu'un morphisme $f : X \to Y$ dans $(\mathrm{Ssimpl})$
est un homotopisme si et seulement si il en est ainsi de $T(f) : T(X) \to
T(Y)$, d'où des foncteurs
\[
S' : (\mathrm{Cat})' \to (\mathrm{Ssimpl})' \quad\text{et}\quad T' :
(\mathrm{Ssimpl})' \to (\mathrm{Cat})'
\]
entre les catégories « homotopiques » construites avec $(\mathrm{Cat})$ resp.\
$(\mathrm{Ssimpl})$, qui sont quasi-inverses l'un de l'autre.

De plus, Quillen construit un isomorphisme canonique et fonctoriel dans
$(\mathrm{Cat})'$ entre $C$ et la catégorie opposée $C^{\circ}$, ou ce qui
revient au même, un isomorphisme canonique et fonctoriel dans
$(\mathrm{Ssimpl})'$ entre $S(C)$ et $S(C^{\circ})$. La définition est telle
que le foncteur induit sur les systêmes locaux sur $C$ transforme le foncteur
contravariant $F$ sur $C$, transformant toute flêche en flêche inversible, en
le foncteur covariant (i.e.\ contravariant sur $C^{\circ}$) ayant mêmes
valeurs sur les objets, et obtenu sur les flêches en remplaçant $F(u)$ par
$F(u)^{-1}$ ; en d'autres termes, l'effet de l'homotopisme de Quillen sur les
groupoïdes fondamentaux est l'isomorphisme évident entre les groupoïdes
fondamentaux de $C$ et de $C^{\circ}$, compte tenu que le deuxième est
l'opposé du

\page{3}

premier. Comme application, Quillen obtient une interprétation faisceautique
de la cohomologie d'un ensemble semi-simplicial à coéfficients dans un systême
local \uncertain{covariant} $F$ (défini classiquement par le complexe cosimplicial des
$C^{n}(F) = \coprod_{x \in X_{n}} F(x)$) : on considère le systême local
contravariant défini par $F$, on l'interprète comme un faisceau sur $T(X)$
i.e.\ objet de $(\mathrm{Ssimpl})/X$, et on prend sa cohomologie.

\marginal{\uncertain{$\prod$} !}\note{en regard de la formule, à la main :
un signe de produit suivi d'un point d'exclamation, qui corrige la somme
tapée. Plus haut, « covariant » : la frappe porte « con », trois lettres biffées à la machine, puis « variant », de sorte que « contravariant » n'est pas exclu ; « local » est tapé en interligne}

--- Cependant, quand $F$ est un systême de coéfficiens covariant pas
nécessairement local, on n'a tjrs pas d'interprétation de ses groupes de
cohomologie classiques en termes faisceautiques ; ni, lorsque $F$ est
contravariant, de son homologie, ou inversement de sa cohomologie faisceautique
en termes classiques.

\marginal{Si !}\note{dans la marge, en regard de ce paragraphe, qu'un grand
crochet manuscrit embrasse}

A propos de la notion de foncteur qui est un homotopisme. Quillen montre qu'un
tel foncteur $F : C \to C'$ induit une équivalence entre la sous catégorie
triangulée $D^{b}_{\mathrm{lc}}(C')$ de la catégorie dérivée bornée de celle
des faisceaux abéliens sur $C'$, dont les faisceaux de cohomologie sont des
systêmes locaux, et la catégorie analogue pour $C$ ; et réciproquement. On
peut dans cet énoncé introduire aussi n'importe quel anneau de base (à
condition de le supposer $\neq 0$ dans le cas de la réciproque) ; la partie
directe vaut aussi avec un anneau de boefficients pas néc.\ constant, mais
constant tordu. Je pense que ce résultat (facile) doit pouvoir se généraliser
ainsi : Soit $f : X \to X'$ un morphisme de topos qui soit tel que pour tout
faisceau loc.\ const.\ sur $X'$, $f$ induise un isomorphisme sur les
cohomologies (avec cas non commutatif inclus). Supposons que \struck{$X$ (resp.\ $X'$)} \add{$X$ et $X'$}
soit\note{« $X$ et $X'$ » est écrit à la main au-dessus de « $X$
(resp.\ $X'$) » biffé ; « soit » et « trivial » restent au singulier} \emph{localement homotopiquement trivial}, i.e.\ que
pour tout entier $n$ \add{$\geq 1$}, tout objet $U$ \struck{de $X'$} ait un recouvrement par des
$U_{i} \to U$, tels que a) tout systême local sur $U$ devient constant sur
$U_{i}$, et toute section sur $U$ devient constante sur $U_{i}$ et b) pour
tout groupe abélien $G$, les $H^{j}(U, G) \to H^{j}(U_{i}, G)$ sont nuls pour
$1 \leq j \leq n$. Alors le foncteur $D^{b}_{\mathrm{lc}}(X') \to
D^{b}_{\mathrm{lc}}(X)$ induit par $f$ est

\marginal{Attention, cette condition n'est \uncertain{typiquement} \emph{pas}
satisfaite par les schémas \uncertain{avec} leur topologie \uncertain{étale} (mais
bien par \struck{\ill{}} \ill{} loc.\ noeth.\ avec top.\
Zariski).}\note{note manuscrite au pied de la page, appelée par un signe de
renvoi porté dans la marge en regard de la définition, que longe un double
trait vertical ; la flèche de « $f : X \to X'$ » manque à la frappe}

\page{4}

une équivalence). Même énoncé si on met dans le coup un systême local
d'anneaux sur $X'$. Enfin, $f$ induit une équivalence entre la catégorie des
coéfficients locaux sur $C$ et celle des coéff.\ locaux sur
$C'$.\note{la phrase a été retouchée à la machine, en plusieurs couches
superposées : « pleinement fidèle (resp.\ » est biffé en tête de page, et
« foncteur pl.\ fid.\ (resp.\ une équivalence) » surcharge « isomorphisme ».
Suit, biffé à la machine : « (?? peut-être cette propriété doit être mise dans
les hypothèses préliminaires sur $f$ ??) »} Principe de démonstration : on
commence par prouver ce dernier résultat, en notant que si un topos est
localement hom.\ trivial, il est loc.\ connexe et loc.\ simplement connexe,
d'où une bonne théorie du $\pi_{0}$ et du $\pi_{1}$ (qui sont ici discrets), et
on est ramené à un cas particulier du critère d'homotopisme de Artin et Mazur,
savoir un critère cohomologique pour qu'un homomorhisme de groupes $G \to H$
(ici des groupes $\pi_{1}$ de $X$, $X'$) soit un isomorphisme : il doit
induire des isomorphismes sur les $H^{0}$ et $H^{1}$ (y inclus dans le cas non
commutatif ...). On prouve la pleine fidélité en se ramenant par la suite
spectrale aux cas des $\mathrm{Ext}^{i}$ de coéfficients locaux. Par la suite
spectrale encore, cela résultera du fait suivant : si $X$ est localement hom.\
trivial, alors la catégorie des faisceaux abéliens loc.\ constants est
\emph{stable par $\mathrm{Ext}^{i}$}, et le foncteur $M \mapsto M_{X}$ des
$(\mathrm{Ab})$ dans $X_{\mathrm{ab}}$ commute aux dits $\mathrm{Ext}^{i}$. En
fait, $X$ et $X'$ étant loc.\ hom.\ triviaux, les conditions suivantes sur $f$
seront équivalentes :
\begin{itemize}
\item[a)] $f$ est un homotopisme, i.e.\ induit pour tout systême local (pas
néc commutatif) sur $X'$ un isomorphisme sur les $H^{i}$.
\item[b)] $f$ induit une équivalence $D^{b}_{\mathrm{lc}}(X') \to
D^{b}_{\mathrm{lc}}(X)$.
\item[a')] $f$ induit une équivalence entre la catégorie des systêmes locaux
abéliens sur $X'$ et $X$, et des isomorphismes sur les $H^{i}$ correspondants
(donc on ne prend ici que des coéfficients commutatifs).
\end{itemize}

\marginal{?}\note{dans la marge, en regard de la phrase qui suit}

J'ignore si on peut dans a) se borner aux systêmes locaux commutatifs.

L'équivalence entre a) et b) fournit une première justification ou motivation
pour définir des types d'homotopie via la catégorie
$D^{b}_{\mathrm{lc}}(X)$,\note{l'exposant et l'indice de ce
dernier $D$ sont ajoutés à la main}

\page{5}

éventuellement muni de la sous-catégorie pleine de tous les systêmes locaux sur
$X$, et du foncteur cohomologique sur $D^{b}_{\mathrm{lc}}(X)$ à valeurs dans
ladite catégorie, et bien sûr du produit tensoriel (mais alors on sort de
$D^{b}$ pour entrer dans $D^{-}$, redactor demerdetur).

\subsection*{II \quad $n$-catégories, catégories $n$-uples, et
(Gr)-catégories.}

\note{sur cette page et la suivante, la main barre les exposants tapés $0$,
$1$, $00$, ... et les récrit en indices ; on donne la forme corrigée}

Appelons bi-catégorie une catégorie dans $(\mathrm{Cat})$, i.e.\ un triple
$(C_{0}, C_{1}, p_{1}, p_{2}, \pi)$, où $C_{0}$ est une catégorie, $p_{1},
p_{2} : C_{1} \rightrightarrows C_{0}$ deux foncteurs, $\pi : (C_{1}, p_{1})
\times_{C_{0}} (C_{1}, p_{2}) \to C_{1}$ un foncteur, avec les axiomes
habituels. Interprétant $(\mathrm{Cat})$ comme la sous-catégorie pleine de
$(\mathrm{Ssimpl})$ des objets qui ..., et de même pour les catégories à
valeurs dans une catégorie telle $(\mathrm{Cat})$, on trouve que la catégorie
des $(\mathrm{Cat})$-catégories $(\mathrm{Bicat})$ est équivalente à la
sous-catégorie pleine de la catégorie des ensembles bi-semisimpliciaux
\[
\begin{array}{cccc}
X_{00} & X_{01} & \cdots & X_{0n} \\
X_{10} & X_{11} & \cdots & X_{1n} \\
\vdots & & & \vdots
\end{array}
\]
formée des objets qui ligne par ligne et colonne par colonne satisfont la
condition d'exactitude familière (transformer sommes amalgamées en produits
fibrés). De même la catégorie $(\mathrm{Tricat})$ des
$(\mathrm{Bicat})$-catégories s'interprète en termes d'objets
tri-semisimpliciaux, etc.

Exemple : soit $C$ une catégorie, posons $C_{0} = C$, $C_{1} =
\mathrm{Fl}(C)$, avec structures $p_{1}$, $p_{2}$, évidentes (d'où une
$(\mathrm{Cat})$-catétorie dont les composantes semi-simpliciales supérieures
sont les catégories $\mathrm{Fl}^{n}(C)$ des $n$-flêches de $C$ ...).
L'ensemble bi-semisimplicial correspondant a pour composantes les ensembles
$\mathrm{Hom}(\Delta^{n} \times \Delta^{m}, C)$. Par exemple, les objets de
$X_{11}$ sont les carrés commutatifs de flêches de $C$, qu'on peut interpréter
(de deux manières différentes) comme

\page{6}

morphismes de flêches. Notons une symmétrie évidente dans la catégorie
$(\mathrm{Bicat})$, interprétée en termes d'ensemble bi-semisimpliciaux,
obtenue en échangeant les lignes et les colonnes dans de tels ensembles bi
... Un ensemble bisemisimplicial qui ... peut donc de deux façons
différentes être considéré comme définissant une bicatégorie. De même le
groupe symétrique $\mathrm{Sym}_{n}$ opère sur la catégorie
$(n\text{-upcat})$ des catégories $n$-uples ...

Une 2-catégorie peut être définie comme constituée par un ensemble $C_{0}$
(l'ensemble des objets), qu'on interprètera comme une catégorie discrète, plus
la donnée pour tout $x, y \in C_{0}$ d'une catégorie
$\underline{\mathrm{Hom}}(x, y)$, plus des foncteurs composition entre celles-ci
... Prenant la catégorie somme $C_{1}$ des $\mathrm{Hom}(x, y)$, On peut
exprimer la donnée de ces foncteurs de composition, et des sources et buts
$x$, $y$ dans les catégories $\mathrm{Hom}(x, y)$, comme définis par des
foncteurs « source » et « but » $C_{1} \to C_{0}$, et un foncteur de
composition $C_{1} \times_{C_{0}} C_{1} \to C_{1}$, les axiomes des
2-catégories correspondant exactement aux axiomes des
$(\mathrm{Cat})$-catégories. De cette façon, la catégorie des 2-catégories
est équivalente à la catégorie des $(\mathrm{Cat})$-catégories qui ont une
« catégorie d'objets » qui est discrète, ou encore à la catégorie des
Bicatégories qui, interprétées comme ensembles bi-semisimpliciaux, ont comme
première ligne un ensemble semisimplicial discret (i.e.\ un foncteur constant
sur $\Delta$). Il parait qu'on a une interprétation simple dans le même
esprit de la notion de $n$-catégorie en termes d'ensembles
$n$-semi-simpliciaux. Demander à Quillen relations entre la notion de
$m$-Catégorie dans la catégorie $(n\text{-Cat})$, et celle de ensemble
$(m+n)$-semisimplicial.\note{un double trait vertical manuscrit, dans la
marge, en regard de ces deux phrases. Suit, biffé à la machine : « En
particulier entre cette dernière notion et celle de »} L'interprétation
standard semisimpliciale permet d'interpréter

\page{7}

\marginal{Attention à la catégorie $\underline{\mathrm{Hom}}_{\mathrm{grab}}$, qui
s'identifie à la catégorie définie par un $\mathrm{RHom}_{\mathbf{Z}}(C_{\bullet},
C'_{\bullet})$, ce point devant être élucidé !}\note{note manuscrite
encadrée, en tête de la page, appelée par un signe de renvoi}

de même les $(\mathrm{Cat})$-groupes (ou groupes à valeurs dans
$(\mathrm{Cat})$) comme étant les groupes semisimpliciaux qui « sont » des
catégories i.e.\ qui satisfont la condition d'exactitude familière, permettant
d'en déterminer tous les composants à l'aide de $C_{0}$ et $C_{1}$. Appliquant
le tapis de Dold, dans le cas commutatif, on trouve que les
$(\mathrm{Cat})$-groupes abéliens forment une catégorie équivalente à celle des
complexes de chaines de groupes abéliens qui sont de longueur 1 :
\[
0 \longleftarrow C_{0} \longleftarrow C_{1} \longleftarrow 0 .
\]
Quand on explicite la $(\mathrm{Cat})$-groupe abélien associé à un tel
complexe de chaines, on trouve « ma » construction chérie de la « catégorie de
Picard » associée au complexe : objets ceux de $C_{0}$, flêches de $x$ dans
$y$ les $z$ dans $C_{1}$ tels que $dz = y - x$. On trouve donc à isomorphisme
près tous les groupes abéliens de $(\mathrm{Cat})$ ainsi. Un complément
important est que toute catégorie \struck{(resp.\ et} de Picard
\struck{commutative} est équivalente (au sens technique qu'on ne précise pas
ici) à un vrai $(\mathrm{Cat})$-groupe \struck{(resp.\ $(\mathrm{Cat})$-groupe
commutatif)}.\note{les deux « resp. » et « commutative », « commutatif) » sont
tapés en interligne et biffés à la main}

\marginal{Quand on \uncertain{simplifie} mod « \uncertain{Grab} »-équivalence, cela correspond à
travailler dans la catégorie \emph{dérivée} des complexes, \uncertain{seuls invariants} \ill{}
\uncertain{$H_{0}$, $H_{1}$} (\uncertain{ou} $\pi_{0}$, $\pi_{1}$)).}\note{à la main, à la suite de la ligne tapée qui finit sur « $(\mathrm{Cat})$-groupe », puis sur deux lignes serrées entre
les lignes tapées, qui repartent de la marge gauche. La copie à l'encre bleue (dossier 162-5, p. 38) lit clairement « Quand on simplifie mod « Grab »-équivalence » et « seuls invariants sont $H_{0}$, $H_{1}$ » ; la photocopie ne les porte que pâlis. Plus haut dans la marge,
en regard des lignes sur la catégorie de Picard, une note manuscrite écrite de
biais est barrée de traits obliques parallèles ; elle ne se lit pas}

On notera qu'on a une classification beaucoup plus restreinte que pour les
H-espaces (resp.\ les H-espaces homotopiquement comm.)\ mod équivalence
d'homotopie, de façon précise entre les groupes resp.\ les groupes abéliens de
la catégorie homotopique. La raison en est qu'on est beaucoup plus restrictif
dans le contexte catégorique pour le sens de l'associativité (resp.\ et de la
commutativité), puisqu'on ne demande pas seulement l'existence d'isomorphismes
d'associativité (qui suffirait pour avoir un groupe dans la catégorie
homotopique $(\mathrm{Cat})' \simeq (\mathrm{Ssimpl})'$, modulo l'existence
d'unités et d'inverses) mais qu'on suppose \add{\emph{donnés}} ces
isomorphismes, de façon à satisfaire à certaines compatibilités : c'est ce qui
permet, me semble-t-il, de prouver qu'à équivalence près, on peut supposer
qu'il s'agit de vrais groupes dans $(\mathrm{Cat})$\struck{, resp.\ de vraie
commutativité}.\note{« donnés » est tapé et souligné à la main ; une croix
manuscrite de renvoi après « d'associativité », et un crochet manuscrit
autour des deux dernières lignes}

\page{8}

On doit trouver une approximation meilleure d'un cran à la classification
topologique en prenant des $(2\text{-Cat})$-groupes resp.\ des
$(2\text{-Cat})$-groupes abéliens, ou plus généralement des
« Gr-2-catégories » resp.\ « Grab-2-catégories ». (NB Une Grab-catégorie
définit une Grab-2catégorie comme une catégorie définit une 2-catégorie ; du
point de vue semi-simplicial, cette opération serait analogue à celle de
« cosquélette${}^{\times}$ tronqué » ; Mais également une Grab-2-catégorie
définit une Grab-catégorie${}^{\times}$, savoir la catégorie des
automorphismes de l'objet unité ; cette opération serait analogue à celle de
« tuage des groupes d'homotopie en dim $> 1$ ».) Pour la notion grossière (Gr
ou Grab) l'idée de ma lettre à Verdier était qu'on doit trouver la même chose
que les tronqués au cran 2, dans la \emph{catégorie homotopique} des
H-espaces --- i.e.\ ce qu'on déduit des H-espaces par tuage des groupes
d'homotopie en dimensions $> 2$.\note{les signes $>$ sont récrits à la main}

(Demander à Quillen relations exactes entre objets-groupes et objets connexes
de la catégorie homotopique ponctuée : forment-ils des catégories équivalentes
par le foncteur espace des lacets\supplied{ ?}).\note{un double trait vertical
manuscrit en marge de la parenthèse ; « ponctuée » est tapé en interligne et mis à sa place, après « homotopique », par un trait}

\marginal{non, on a un foncteur pl.\ fidèle (?) mais pas ess.\
surjectif ....}\note{note écrite de biais dans la
marge gauche, en regard de la question}

Pour les \emph{vrais} groupes abéliens de $(2\text{-Cat})$, l'approche de
Quillen consiste à les identifier d'abord à des groupes abéliens de la
catégorie $(\mathrm{Bicat})$, donc à des groupes abéliens bi-semisimpliciaux
particuliers (conditions sur chaque ligne et chaque colonne). Appliquant la
version bi-semisimpliciale du tapis de Dold, on trouve que la catégorie de
ces objets est équivalente à la catégorie des bicomplexes de chaines de
groupes abéliens qui sont nuls en dehors du carré $0 \leq i, j \leq 1$. La
\struck{colonne} \add{cond.}\ que la première ligne corresponde à un groupe
semisimplicial trivial signifie alors qu'elle n'a que le terme de degré zéro,
i.e.\ que le bicomplexe est réduit à

\begin{tikzcd}
C_{00} & 0 \arrow[l] & 0 \arrow[l] \\
C_{10} \arrow[u] & C_{11} \arrow[l] \arrow[u] & 0 \arrow[l]
\end{tikzcd}

\note{les deux « $\leftarrow 0$ » de droite sont ajoutés à la main}

\page{9}

Ces bicomplexes s'identifient aux complexes de chaines de longueur 2. Donc la
catégorie de $(2\text{-Cat})$-groupes abéliens est équivelente à celle des
complexes de chaines de longueur 2. L'explicitation de la 2-catégorie associée
à un tel complexe de chaines est mon interprétation chérie. Il faudrait encore
expliciter 2-Grab-catégorie des homomorphismes (en un sens à préciser) entre
deux telles 2-Grab-catégories strictes, en termes d'un complexe
$\mathrm{RHom}$.

L'idée qui se dégage de cette approche, et que Quillen on espère va expliciter,
c'est que la notions de groupe abélien dans $(n\text{-Cat})$ correspond à
celle de complexe de chaines de longueur $n$, en raisonnant soit mod.\
isomorphisme de part et d'autre, soit mod.\ une notion naturelle de
$n$-équivalence d'un côté, et dans la catégorie dérivée de l'autre. Ceci doit
pouvoir se préciser en utilisant des $\mathrm{Hom}^{\bullet}$ resp.\ des
$\mathrm{RHom}^{\bullet}$.

Si maintenant $C$ est une catégorie avec un « produit tensoriel » associatif
(et éventuellement unitaire), i.e.\ un bifoncteur avec des isomorphismes
d'associativité (mais sans inverse ...), on peut lui associer sans doute
une Gr-catégorie universelle (tout comme à un monoïde on associe un groupe
symmétrisé), d'où des invariants, qu'on pourra désigner par $K_{0}(C)$ et
$K_{1}(C)$. On suppose qu'on pourra de même trouver, pour tout $n$, une
Gr-$n$-catégorie universelle où envoyer $C$ de façon compatible avec $\otimes$
(tout cela dans un sens à préciser), dont les invariants $\pi_{n}$ seraient
par définition les $K_{n}(C)$. Un cas intéressant est celui où $C$ est une
catégorie avec produits finis (par exemple une catégorie additive) et
$\otimes$ est le produit ; Mais on notera que dans le cas d'une catégorie non
seulement additive mais abélienne, il y a une notion plus adéquate du $K_{0}$,
qui devrait correspondre à une notion plus adéquate de la suite des $K_{i}$,
inspirée du tapis de Roos ; sans doute c'est en termes de (« vraie »)
catégories triangulées\note{un crochet ouvrant manuscrit devant « Mais on
notera »}

\page{10}

qu'il conviendra de poser les définitions adéquates). --- Peut-être Quillen
a-t-il une définition semi-simpliciale plus directe, sans passer par des
$n$-catégories, des $K_{i}$ et plus précisément du H-espace qui devrait
correspondre à $C$.

\marginal{oui}\note{dans la marge, avec un double trait vertical, en regard
de la phrase qui précède}

Ce n'est certainement pas en général l'ensemble semi-simplicial $SC$ associé à
$C$, puisque dans le cas d'une catégorie additive, $C$ est homotopiquement
trivial (en effet le foncteur identique est « homotope » au foncteur nul,
puisqu'il suffit d'avoir un homomorphisme d'un foncteur dans l'autre ; d'où une
équivalence d'homotopie entre $C$ et la catégorie nulle). Pourtant, il semble
bien que dans le cas où $C$ est une Gr-catégorie i.e.\ admet des inverses,
alors ce soit bien la construction « naïve » $SC$ qui est la bonne, puisqu'on
trouve bien un $\pi_{0}$ et un $\pi_{1}$ qui sont aussi les $\pi_{0}$ et
$\pi_{1}$ au sens catégorie sans Gr-structure.

\marginal{A expliquer : pour des Gr-$n$-catégories, les $\pi_{i}$ au sens
semi-simplicial sont aussi les $\pi_{i}$ au sens Gr-$n$-catégories, savoir les
objets à isom.\ près dans divers $\mathrm{Hom}(1, 1)$ ...\
\struck{Il y a aussi une interprétation de l'opération \uncertain{$\Omega$ :} \ill{} \uncertain{des lacets} en termes du
foncteur d'opération \uncertain{du} shift \uncertain{géométrique} \uncertain{évidente} sur les
$n$-catégories ....}}\note{paragraphe manuscrit sous la frappe, appelé
par un signe de renvoi en marge ; sa seconde moitié est barrée de traits
obliques, sous lesquels elle reste lisible par endroits ; la copie à l'encre bleue (dossier 162-5, p. 41) porte la même phrase}

\page{11}

\subsection*{III \quad Point de vue « motivique » en théorie du cobordisme.}

Soit $C$ la catégorie des variétés différentiables \add{(compactes ?),} (pas
nécessairement orientables), les morphismes étant les classes d'homotopie
d'applications continues. Si $B$ est une catégorie, on s'intéresse aux couples
$(F_{\cdot}, F^{\bullet})$ d'un foncteur covariant et d'un foncteur
contravariant de $C$ dans $B$, satisfaisant les conditions que pour tout $X
\in \mathrm{Ob}\, C$, on a $F_{\cdot}(X) = F^{\bullet}(X)$, et que si on a un
produit fibré ordinaire

\begin{tikzcd}
X' \arrow[r, "g'"] \arrow[d, "f'"'] & X \arrow[d, "f"] \\
Y' \arrow[r, "g"'] & Y
\end{tikzcd}

avec $f$ et $g$ transversaux, alors on a $g^{\bullet} f_{\cdot} = f'_{\cdot}
g'^{\bullet}$ (on pose $f_{\cdot} = F_{\cdot}(f)$, $f^{\bullet} =
F^{\bullet}(f)$). La question est de trouver une $B$ et un couple
$(F_{\cdot}, F^{\bullet})$ qui soient « universels » dans un sens évident.
Quillen prouve que la construction suivante donne une telle solution.

Pour toute variété $X \in \mathrm{ob}\, C$, on définit un anneau commutatif
$B(X)$ (l'anneau de bordisme de $X$) en prenant comme éléments les classes
d'éléments de $C$ au dessus de $X$, pour la relation de cobordisme : $f' : X'
\to X$ et $f'' : X'' \to X$ sont équivalents, s'il existe une variété à bord
$Z$, et un $f : Z \to X$, tel que $\mathrm{bord}\, Z = X' \sqcup X''$ et $f'$
et $f''$ sont induits par $f$. La somme de $B(X)$ est la somme disjointe
(c'est une loi de groupe où tout élément est d'ordre 2), le produit est donné
par le produit fibré, en utilisant le théorème de transversalité de Thom pour
bouger un peu $f''$, de sorte à le rendre transversal à $f'$.

\marginal{\emph{NB} L'anneau est de plus \emph{\uncertain{gradué}} par la
dimension relative ...}\note{en interligne, sous le paragraphe}

Si on a $f : X \to Y$, on définit $f_{\cdot} : B(X) \to B(Y)$ de façon
évidente (toute variété sur $X$ est sur $Y$), et $f^{\bullet} : B(Y) \to B(X)$
par produit fibré (utilisant encore le th.\ de transversalité). De cette
façon, $B(X)$ devient un foncteur-anneau \add{(gradué)} contravariant en $X$,
et un foncteur groupe covariant en

\page{12}

$X$. De plus, on a la formule \add{(du type $g^{\bullet} f_{\cdot} =
f'_{\cdot} g'^{\bullet}$ ci-dessus, impliquant la formule)} de projection
habituelle. Ceci posé, les constructions bien connues en théorie des
correspondances permettent de définir une catégorie $\mathbf{B}$ dont les
objets sont ceux de $C$, et les Hom étant donnés par
\[
\mathrm{Hom}_{\mathbf{B}}(X, Y) = B(X \times Y) .
\]
On a un isomorphisme canonique $\mathbf{B} \simeq \mathbf{B}^{\circ}$, qui
est l'identité sur les objets et la symmétrie $B(X \times Y) \simeq B(Y
\times X)$ sur les flêches, et un foncteur évident $F_{\cdot} : C \to
\mathbf{B}$, qui est l'identité sur les objets et donné par les graphes sur
les flêches. Composé avec la symmétrie précédente, on trouve un foncteur
$F^{\bullet} : C \to \mathbf{B}$.\note{la frappe porte ici $C$, sans le $^{\circ}$ qu'appelle un foncteur contravariant} On vérifie enfin que le couple
$(F_{\cdot}, F^{\bullet})$ satisfait aux conditions envisagées dans l'alinéa 1
(ceci aussi devrait être formalisé au cas des correspondances générales
...). On vérifie enfin qu'il a la propriété universelle qu'il faut.
Lorsqu'on se donne un couple $(F'_{\cdot}, F'^{\bullet})$ à valeurs dans une
catégorie $B'$, il faut trouver une factorisation par $\mathbf{B}$. Elle est
toute trouvée sur les objets, il faut la définir sur les flêches, donc
définir, pour $X$, $Y$ dans $C$, une application
\[
B(X \times Y) \longrightarrow \mathrm{Hom}(X', Y') . \tag{$*$}
\]
Mais si $Z$ est au-dessus de $X \times Y$, d'où $p : Z \to X$ et $q : Z \to
Y$, on lui associe un élément du deuxième membre de $(*)$ comme le composé
\[
X' \xrightarrow{F'^{\bullet}(p)} Z' \xrightarrow{F'_{\cdot}(q)} Y' .
\]
Il faut prouver que l'élément obtenu en remplaçant $Z$ par un $Z_{1}$
cobordant au dessus de $X$ \add{est le même}. Or si $g : T \to X \times Y$
réalise le cobordisme, on voit (regardant seulement $T$ et la décomposition
$\mathrm{bord}(T) = Z \sqcup Z_{1}$) qu'il existe une variété $\overline{T}$
sur $S^{1}$ (obtenue en « doublant » $T$) transversale au dessus d'un couple de
points $s_{0}$, $s_{1}$ de $S^{1}$, tels que $Z_{i}$ soit l'image inverse
\add{dans $\overline{T}$} de $s_{i}$ ($i = 0, 1$) ; d'autre part,

\page{13}

$g$ manifestement se prolonge en $\overline{g} : \overline{T} \to X \times
Y$. Cela nous ramène à prouver que le composé
\[
\overline{T}' \xrightarrow{\alpha_{i}} Z'_{i}
\xrightarrow{\alpha_{i*}} \overline{T}'
\]
ne dépend pas de $i = 0, 1$. Ecrivons $X$ au lieu de $\overline{T}'$, et
considérons le diagramme cartésien d'inclusions

\begin{tikzcd}
X_{s} \arrow[r, "\alpha_{s}"] \arrow[d, "\alpha_{s}"'] & X \arrow[d, "{\Gamma = (p, \mathrm{id}_{X})}"] \\
X \arrow[r, "i_{s}"'] & S^{1} \times X
\end{tikzcd}

on a alors $\alpha_{s}^{*} \alpha_{s*} = \Gamma^{*} i_{s*}$, qui ne dépend
pas de $s$ puisque $i_{s}$ n'en dépend pas à homotopie près. --- Le reste est
sans doute l'AQT.\note{les $\alpha$, les astérisques et les étiquettes du
carré sont portés à la main sur des blancs de la frappe. La première étiquette du composé se lit $\alpha_{i}$, sans astérisque visible, ici comme sur la copie à l'encre bleue (dossier 162-5, p. 45). L'étiquette de droite du carré, pâle sur la photocopie, se lit $\Gamma = (p, \mathrm{id}_{X})$, le signe $=$ souligné}

Sur la strcture des $B_{\cdot}(X)$. Ce sont tous des algèbres graduées sur
$B_{\cdot}(\mathrm{pt}) = \Lambda$, qui est une algèbre graduée à degrés
positifs, réduite à $\mathbf{F}_{2}$ en degré 0. C'est d'ailleurs une algèbre
de polynômes \add{(à une infinité de générateurs $x_{i}$ ($i$
\uncertain{entier $\geq 1$}, \uncertain{$\neq$ de la forme $2^{\nu} - 1$}))} connue ....\note{l'insertion manuscrite, pâle et rognée au bord droit, est lue avec la copie à l'encre bleue (dossier 162-5, p. 45), qui porte « ($i$ entier $\geq 1$, $\neq$ de la forme $2^{\nu} - 1$) »}

\marginal{NB $H^{*}(X)$ $(= H^{*}(X, \mathbf{F}_{2}))$}\note{dans la marge gauche}

On a un homomorphisme évident $B_{\cdot}(X) \to H^{\bullet}(X)$, nul en degrés
$> 0$, qui à tout $Z$ sur $X$ associe l'image de $1 \in H^{0}(Z)$ par
l'homomorphisme de Gysin associé à $Z \to X$. Cet homomorphisme par Thom est
surjectif (toute classe de cohomologie provient d'une image de variété
compacte ...) ;

\marginal{\uncertain{en fait} $H^{\bullet}(X) \simeq B_{\cdot}(X) /
\Lambda^{+} B_{\cdot}(X)$}\note{dans la marge gauche}

On ppouve même qu'il existe un relèvement fonctoriel en $X$
\[
H^{\bullet}(X) \longrightarrow B_{\cdot}(X),
\]
\add{(il est douteux qu'il soit} compatible avec multiplication\supplied{)} tel que l'on ait :\note{la main met « compatible avec
multiplication » entre parenthèses et écrit au-dessus « il est douteux qu'il
soit » ; un point d'interrogation dans la marge}

\struck{a)} L'homomorphisme correspondant $H^{\bullet}(X)
\otimes_{\mathbf{F}_{2}} \Lambda \simeq B_{\cdot}(X)$ est un isomorphisme (ce
sera vrai pour tout relèvement individuellement ...), et

b) \add{[} \add{Quillen ignore si} \struck{Le} relèvement est compatible aussi avec les hom.\ de Gysin.\note{« a) » est raturé à la main ; devant b), un crochet manuscrit ouvre la clause, et « Quillen ignore si », écrit au-dessus de « Le » biffé, en fait une question. Un grand trait en arc dans la marge
en regard de a) et b)}

Ainsi, la catégorie $\mathbf{B}$ admet un foncteur pleinement fidèle dans la
catégorie des modules gradués libres de type fini sur $\Lambda$, les modules de
type fini qu'on obtient ainsi ayant des générateurs à degrés $\leq 0$ et
satisfaisant (pour $X$ connexe i.e.\ quand il y a un seul générateur de degré
zéro) la dualité de Poincaré. Quand on prend la catégorie \emph{quasi-abélienne}
\add{$\widetilde{\mathbf{B}}$} associée à $\mathbf{B}$, on trouve donc tous les modules gradués libres de type
fini à générateurs en degrés $\leq 0$ seulement. C'est une

\marginal{i.e.\ l'enveloppe « Karoubienne »}\note{de biais dans la marge
gauche, en regard de « quasi-abélienne », que la main souligne d'un trait
ondulé}

\page{14}

catégorie additive (et même quasi-abélienne) avec $\otimes$ qui a toutes les
vertus, sauf d'être abélienne, et sauf d'admettre des
$\underline{\mathrm{Hom}}$ internes.

\marginal{= karoubienne + additive}\note{en tête de la page,
au-dessus de « quasi-abélienne », qui est tapé en interligne}

Pour trouver ces derniers, on peut si on veut paraphraser la construction de
la catégorie des motifs, \add{en} considérant l'objet \add{$\Lambda(-1)$} de
$\mathbf{B}$ conoyau de $p^{*} : e \to S$, où $p : S^{1} \to e$ est
l'application du cercle sur un point. C'est un objet de poids 1 (le poids est
défini par le degré des générateurs, changés de signe). On constate aussitôt
que le foncteur $M \mapsto M(-1) = M \otimes \Lambda(-1)$ de
$\widetilde{\mathbf{B}}$ dans lui-même est pleinement fidèle, ce qui permet
alors, par une construction formelle standard, de compléter
$\widetilde{\mathbf{B}}$ en une catégorie où $\Lambda(-1)$ devient inversible.
Il se trouve alors que dans cette dernière les $\underline{\mathrm{Hom}}(X, Y)$
internes existent et sont isomorphes à $\check{X} \otimes Y$ ... en fait,
on trouve maintenant la catégorie de \emph{tous} les modules gradués libres de
type fini sur $\Lambda$.\note{les tildes sur $\mathbf{B}$ et le signe sur
$\check{X}$ sont portés à la main ; « en » remplace une frappe biffée à la
machine ; plus bas, « de \emph{tous} » porte une lettre biffée à la machine après « de », et « convient » est une frappe surchargée qu'on lirait aussi « convinmt »}

Bien sûr, cette catégorie n'est pas encore abélienne pour autant, puisque
$\Lambda$ n'est pas un corps ; si on veut définir la catégorie abélienne
engendrée, on s'attend à trouver la catégorie des modules de \struck{type}
\add{présentation} finie sur $\Lambda$. (\add{qui est abélienne bien que
$\Lambda$ ne soit pas noethérien, car $\Lambda$ est} \struck{est-il
noethérien, i.e.\ le nombre de générateurs est-il fini ?} \add{lim
d'anneaux noeth.\ à homs de transition \emph{plats}}).\note{« lim », à la main, sans flèche visible sous le mot}

Cette construction ressemblerait à celle d'une théorie des motifs où on
n'aurait pas passé à l'équivalence homologique pour les cycles, mais où on
aurait gardé l'équivalence rationnelle, --- qui est bien trop fine à beaucoup
d'égards. Pour trouver un équivalent plus étroit, il \uncertain{convient} de reprendre la
théorie prédédente, on pourrait formuler le même Pb universel que plus haut,
mais en prenant l'équivalence homologique (via le graphe) des morphismes de
variétés au lieu de l'équivalence homotopique. \struck{Mais alors il semble
que l'on trouve comme catégorie solution celle dont les objets sont les
varié-}

\page{15}

\struck{tés, les morphismes \emph{toutes} les correspondances cohomologiques
(en vertu du thorème de Thom rappelé plus haut). On trouve donc un foncteur
pleinement fidèle de cette catégorie dans celle des vectoriels gradués
positivement de dim.\ finie sur $\mathbf{F}_{2}$, qui est déjà abélienne, et
si on veut rendre comme plus haut $\mathbf{F}_{2}(-1)$ inversible, on trouve
tous les vectoriels gradués de dim.\ finie sur $\mathbf{F}_{2}$. C'est un peu
trop trivial comme structure, et on peut faire mieux,}\note{le passage, qui
commence à la dernière ligne de la page 14, est biffé à la main : d'un trait
par ligne sur la page 14, de grands traits en zigzag sur la page 15}

Pour trouver une catégorie universelle, on associe à toute $X$ $S(X)$
l'ensemble des sous-variétés de $X$, modulo équivalence homologique. Donc
$S(X) \subset H^{\bullet}(X)$, et on prouve qu'on trouve en fait un sous-anneau
de $H^{\bullet}(X)$ (grâce à un théorème de Thom qui caractérise $S(X)$ en
termes d'opérations de Steenrod, cf plus bas). On prouve de même (par le même
théorème) que les $S(X \times Y)$ sont stables par composition de
correspondances cohomologiques (cela ne semble pas trivial, puisque les
$S(X)$ ne sont pas stables par Gysin ...). D'où une catégorie additive dont
les objets sont ceux de $C$, les flêches sont les \add{él.\ des} $S(X \times
Y)$, Il semble qu'elle doive jouer le rôle universel que j'ai envisagé plus
haut. En tous cas, sa définition géométrique est fort simple Nous allons en
donner une interpétation algébrique. Pour ceci, considérons l'anneau des
opérations cohomologiques « stables ». C'est un anneau gradué, avec
application diagonale anticommutative, i.e.\ commutative (car.\ 2 !),
associative et unitaire. Son dual gradué est une algèbre graduée
anticommutative donc commutative, avec application diagonale associative et
unitaire ; il est réduit à $\mathbf{F}_{2}$ en degré zéro. Donc il corespond à
un schéma en groupes affine $G$ sur $\mathbf{F}_{2}$, qui est d'ailleurs le
foncteur des $\otimes$-automorphismes additifs du foncteur $H^{\bullet}$ sur
$C$ \add{qui opèrent trivialement sur \struck{\ill{}} $H^{1}(S^{1})$.} \struck{\ill{}}\note{après
« stables », tapé en interligne, une définition tapée est biffée à la main :
« des endomorphismes du foncteur $X \mapsto H^{\bullet}(X)$ qui sont additifs,
et qui en degré zéro se réduisent à un multiple de l'identité ». Un long trait
vertical et un signe de renvoi en marge, en regard de ce passage ; l'addition
manuscrite finale biffe un mot avant $H^{1}(S^{1})$ et se termine sur un mot raturé}

\page{16}

Le théorème de Thom est alors qu'un $x \in H^{\bullet}(X)$ est dans $S(X)$
sss il est fixe par $G$, i.e.\ satisfait $A^{+} \cdot x = 0$. (Pourquoi cette
condition est-elle déjà nécessaire ?) Une autre façon de le voir est de dire
que le foncteur $H^{\bullet}$ considéré comme allant de la catégorie
$\mathbf{S}$ précédente dans la catégorie des représentations de $G$ dans des
vectoriels sur $\mathbf{F}_{2}$, est pleinement fidèle. Quand on prend
l'enveloppe pseudo-abélienne, on trouve une sous-catégorie pleine de la
catégorie des représentations de $G$, mais il est douteux qu'on trouve une
catégorie abélienne, voire la catégorie de toutes les représentations de $G$ de
dimension finie sur $\mathbf{F}_{2}$. (Quelle est la structure de $G$, qui est
prolisse ; d'après Quillen, il serait pro-unipotent ....) Ici $G$ bien sûr
joue le rôle d'un groupe de Galois motivique, mais contrairement à la situation
familière, il semble qu'il soit fortement non réductif (et d'ailleurs le corps
de définition dans tout ceci est $\mathbf{F}_{2}$, et non $\mathbf{Q}$ !).

\marginal{\uncertain{oui}, car son anneau affine est une \uncertain{algèbre} de polynômes à une
infinité d'indéterminées, \struck{\ill{}} \uncertain{réunion} d'\uncertain{algèbres} \ill{}
à \uncertain{engendrement} \ill{} stables par $\Delta$}\note{note de sept lignes écrite
de biais dans la marge gauche, en regard de la parenthèse sur $G$, qu'une
flèche y rattache ; le premier mot est rogné au bord. La copie à l'encre bleue (dossier 162-5, p. 48) porte la même note, plus lisible}

Quillen\struck{t} donne une généralisation de la construction d'une catégorie
universelle $\mathbf{B}$ au cas où on travaille avec des variétés pas
néeessairement propres (mais $f_{\cdot}$ n'étant supposé défini que pour $f$
propre), et quand on introduit une condition d'orientation \add{stable} sur
les \emph{morphismes} envisagés, via un morphisme $H \to BO(\infty)$, où $H$
est un H-espace (commutatif, associatif et unitaire ?) : une orientation de $f
: X \to Y$ sera un relèvement homotopique de $X \to BO$, défini par $T_{X} -
f^{\bullet}(T_{Y})$, en $X \to H$. les éléments de $H(X, Y)$ seront définis par des
diagrammes de morphismes orientés $Z \to X$, $Z \to Y$, le deuxième étant
propre, modulo bords de tels. En tant qu'ensemble, il se définit aussi en
termes de la variété $X \times Y$ munie simplement de la famille des supports
formée des ensembles fermés propres sur $Y$ ...\note{« stable » est tapé
en interligne et rattaché à sa place par un trait manuscrit ; le « t » de « Quillent » est barré à la main}

\marginal{Les $B^{H}_{m}(X)$ ($X$ compact disons) peuvent se décrire
homotopiquement à la Thom, comme les groupes d'homotopie stables d'espaces
$\underline{\mathrm{Hom}}(X, MH(N-m))$ ...}

\marginal{cf exposé Karoubi à Bourbaki ...}\note{deux notes manuscrites au
pied de la dernière page ; la seconde est d'une encre plus pâle}

\end{document}
